\documentclass[12pt]{article}
\usepackage[margin=1in]{geometry}
\usepackage{tikz}
\usetikzlibrary{arrows.meta}
\usepackage[american voltages, european currents]{circuitikz}
\usepackage{amsmath}
\usepackage{amssymb}
\usepackage[table,dvipsnames]{xcolor}
\usepackage{pgfplots}
\pgfplotsset{compat=1.18}
\usepackage{float}
\usepackage{caption}
\usepackage{parskip}

% -- colors from source file --
\definecolor{myblue}{RGB}{0,103,172}
\definecolor{mygray}{RGB}{240,240,240}

\begin{document}

\begin{center}
  {\LARGE\textbf{TikZ / CircuitiKZ Figures}}\\[0.4em]
  {\large\texttt{Lab8Instructions.tex}}
\end{center}
\vspace{1em}
\noindent Edit the TikZ code in this file, then recompile
with \texttt{pdflatex} to preview changes.  Run
\texttt{extract\_tikz\_figures.py} to regenerate the PNGs.

\clearpage

\noindent\textbf{Label:} \texttt{fig:INA125\_complete}\\[0.5em]
\begin{figure}[H]
  \centering
  \begin{circuitikz}
  	\ctikzset{muxdemux/outer label font={\tiny\ttfamily\color{blue}}}
  	\tikzset{myICwl/.style={muxdemux,
  			muxdemux def={Lh=12, Rh=12, w=4,square pins=1,
  				NR=7, NL=6, NB=1, NT=2,},
  			muxdemux label={
  				L1=Vref2.5,
  				L2=VrefOut,
  				L3=Vin+,
  				L4=Rg, 
  				L5=Rg, 
  				L6=Vin-,
  				R1=Vref 10, 
  				R2=Vref 5,
  				R3=Vref BG,
  				R4=Vref COM, 
  				R5=IAref,
  				R6=Sense,
  				R7=Vo,
  				T1=V+,
  				T2=$\overline{\text{Sleep}}$, % ADDED BAR HERE
  				B1=V-,
  				LU1=14, LU2=4,LU3=6, LU4=8, LU5=9,LU6=7,
  				RU1=16, RU2=15, RU3=13, RU4=12,RU5=5,RU6=11,RU7=10,
  				TL1=1,TR2=2,
  				BL1=3,
  				N=INA125}
  			,}
  	}
  	
  	% Main IC Node
  	\draw (7, -2) node[myICwl] (IC) {} ;
  	
  	% PIN 3 CONNECTION (REMAINING AT HALFWAY POINT)
  	\draw (IC.bpin 1) -- ++(0,-0.75) coordinate(cap_junction) 
  	-- ++(0,-0.75) node[anchor=north] {$-5V$};
  	
  	\draw (cap_junction) to[short,*-] ++(-1,0) to[C, l_=$0.1\mu F$] ++(0,-1) node[ground]{};
  	
  	% --- Rest of the circuit ---
  	\draw (6-0.1,0.8)  to[short,-] ++ (-6+0.1,0) to[short,-] ++ (0,-0.5);
  	\draw (6-0.1,-0.32)  to[short,-] ++ (-1+0.1,0) to[short,-*] ++ (0,+1.1);
  	\draw (6-0.1,-1.44)  to[short,-] ++ (-2+0.1,0) to[short,-] ++ (0,-0.55) to[short,-] ++ (-2,0);
  	\draw (6-0.1,-2.56)   
  	to[R] ++(-2,0) to[short,-] coordinate(A) ++ (0,-1.13) to[vR] ++(+2,0) ;
  	\draw(6-1,-4.3)  node {$100\Omega$};
  	\draw(6-1.25,-2.1)  node {$R_f$};
  	\draw (3.89,-3.7) to[short,*-] ++ (0,-0.389) to[short,-] ++ (0.834,0)  ;
  	
  	\draw (6-0.1,-4.8)  to[short,-] ++ (-2+0.1,0) to[short,-] ++ (0,-0.55) to[short,-] ++ (-7,0) to[short,-] ++ (0,3.37) to[short,-] ++ (+1,0);
  	
  	\draw (8.4,-2)  node[ground]{} ;
  	\draw (8.4,-2.95)  node[ground]{} ;  
  	\draw (8.4,-2.95-0.97)  to[short,-] ++ (0.5,0)  to[short,-*] ++ (0,-0.97)  ;
  	
  	\draw (8.4,-2.95-0.97-0.96)  to[short,-] ++ (0.5,0)  to[R, l=$1 k\Omega$] ++ (0,-2) node[ground]{} ;
  	\draw (8.4,-2.95-0.97-0.96)  to[short,-o] ++ (1,0) node[anchor=west]  {$V_{out}$} ;
  	
  	\draw (7-0.56,1.5)  to[short,-*] ++ (0,2) node[anchor=east]  {$+5V$};
  	\draw (7+0.56,1.6)  to[short,-] ++ (0,1) to[short,-*] ++ (-2*0.56,0)  
  	to[short,-] ++ (-2,0) to[C, l=$0.1 \mu F$] ++ (0,-1)  node[ground]{}; 
  	
  	\draw (0,0) to[short, -] ++(0,+0.5) node[anchor=east]  {$V_\text{ref}$} ;
  	\draw (0,0)        
  	to[R, *-*] ++ (-2,-2)  node[anchor=south] {$V_1$} 
  	to[R, -*] ++ (2,-2) ; 
  	\draw (0,0)    
  	to[R, -*] ++ (2,-2)  node[anchor=south] {$V_2$} 
  	to[R, -*] ++ (-2,-2) node[ground]{} ; 
  	\draw(-2, -0.8) node {$R_1 - \delta R$};
  	\draw(2, -0.8) node {$R_1 + \delta R$};
  	\draw(-2, -0.8-2.5) node {$R_1 + \delta R$};
  	\draw(2, -0.8-2.55) node {$R_1 - \delta R$};
  	
  \end{circuitikz}
  \caption{Complete INA125 instrumentation amplifier circuit with load cell interface. The gain is set by the resistance between pins 8 and 9.}
\end{figure}

\end{document}
